\documentclass[12pt,oneside]{amsart}

%\usepackage[T1]{fontenc}
%\usepackage[utf8]{inputenc}
\usepackage[german]{babel}
\usepackage{geometry}                % See geometry.pdf to learn the layout options. There are lots.
\geometry{a4paper}                   % ... or a4paper or a5paper or ... 
\usepackage{color}
\usepackage[table]{xcolor}
\usepackage{graphicx}
\usepackage[mtphrb,mtpcal,mtpfrak]{mtpro2}
\usepackage{amsmath}
\usepackage{amsthm}
\usepackage{enumerate}
\usepackage{tikz}
\usetikzlibrary{calc}
\usetikzlibrary{arrows.meta}
\usepackage[shortalphabetic]{amsrefs}
\usepackage{wrapfig}
\usepackage[textsize=tiny]{todonotes}
\usepackage[many]{tcolorbox}
\usepackage{pgfplots}
\usepackage{booktabs}
\usepgfplotslibrary{patchplots}
\usepgfplotslibrary{fillbetween}

%\usepackage[no-math]{fontspec}
%\setsansfont[Scale=0.92]{Helvetica Neue LT Std 65 Medium}
%\usepackage[scaled=0.92]{helvet}
%\usepackage{times}



\usepackage{mathspec}
\setmathrm{Times LT Std}
\setmainfont[BoldFont=Times LT Std Bold, ItalicFont=Times LT Std Italic, BoldItalicFont=Times LT Std Bold Italic, Ligatures=TeX]{Times LT Std}

%\setmainfont{Baskerville URW}

%\setsansfont[Scale=0.92,BoldFont=Helvetica Neue LT Std 65 Medium,
%    ItalicFont=Helvetica Neue LT Std 56 Italic,
%    BoldItalicFont=Helvetica Neue LT Std 66 Medium Italic,
%    Ligatures=TeX]{Helvetica Neue LT Std 55 Roman}
\setsansfont[BoldFont=Whitney-Semibold.otf,ItalicFont=Whitney-MediumItalic.otf,Ligatures=TeX]{Whitney-Medium.otf}
\setboldmathrm[BoldFont={Times LT Std Bold}]{Times LT Std Bold}

\newcommand{\red}[1]{\textcolor{red}{#1}}
\definecolor{light-gray}{gray}{0.85}

%%% here we define our palette

\definecolor{spacecadet}{HTML}{0D284C}
\definecolor{munsell}{HTML}{008FA8}
\definecolor{banana}{HTML}{FFD932}
\definecolor{cgblue}{HTML}{007CA5}
\definecolor{isabelline}{HTML}{EAEDEA}

%%%

\makeatletter
\def\@secnumfont{\bfseries}
\def\part{\@startsection{part}{0}%
  \z@{\linespacing\@plus\linespacing}{.5\linespacing}%
  {\normalfont\bfseries\raggedright}}
\def\specialsection{\@startsection{section}{1}%
  \z@{\linespacing\@plus\linespacing}{.5\linespacing}%
  {\color{cgblue}\normalfont\centering}}
\def\section{\@startsection{section}{1}%
  \z@{.7\linespacing\@plus\linespacing}{.5\linespacing}%
  {\color{cgblue}\normalfont\sf\bfseries\large\centering}}
\def\subsection{\@startsection{subsection}{2}%
  \z@{.5\linespacing\@plus.7\linespacing}{-.5em}%
  {\color{cgblue}\normalfont\sf\bfseries}}
\def\subsubsection{\@startsection{subsubsection}{3}%
  \z@{.5\linespacing\@plus.7\linespacing}{-.5em}%
  {\color{cgblue}\normalfont\sf\itshape}}
\def\paragraph{\@startsection{paragraph}{4}%
  \z@\z@{-\fontdimen2\font}%
  \color{cgblue}\sf\normalfont}
\def\subparagraph{\@startsection{subparagraph}{5}%
  \z@\z@{-\fontdimen2\font}%
  \color{cgblue}\sf\normalfont}
\makeatother

\DeclareMathOperator{\supp}{supp}
\DeclareMathOperator{\conf}{conf}
\DeclareMathOperator{\tr}{tr}
\DeclareMathOperator{\im}{im}
\DeclareMathOperator{\id}{id}
\DeclareMathOperator{\vspan}{span}
\DeclareMathOperator{\res}{res}
\DeclareMathOperator{\real}{Re}
\DeclareMathOperator{\imag}{Im}
\DeclareMathOperator{\diff}{diff}
\DeclareMathOperator{\homeo}{Homeo}
\DeclareMathOperator{\Mor}{Mor}
\DeclareMathOperator{\Ob}{Ob}
\DeclareMathOperator{\target}{Target}
\DeclareMathOperator{\source}{Source}
\DeclareMathOperator{\aut}{Aut}
\DeclareMathOperator{\Sch}{Sch}

\theoremstyle{plain}% default
\newtheorem{theorem}{Theorem}[section]
\newtheorem{lemma}[theorem]{Lemma}
\newtheorem{proposition}[theorem]{Proposition}
\newtheorem*{corollary}{Corollary}

\theoremstyle{definition}
\newtheorem{definition}{\color{cgblue}Definition}[section]
\newtheorem{prototypedefinition}{Prototype Definition}[section]
\newtheorem{physicalintuition}{Physical intuition}[section]
\newtheorem{conjecture}{Conjecture}[section]
\newtheorem{beispiel}{\color{cgblue}Beispiel}[section]
\newtheorem{übung}{\color{cgblue}Übung}[section]
\newtheorem{exercise}{\color{cgblue}Exercise}[section]

\theoremstyle{remark}
\newtheorem*{remark}{Remark}
\newtheorem*{note}{Note}

\newenvironment{lösung}
  {\begin{proof}[\color{cgblue}Lösung]\color{spacecadet}}
  {\end{proof}}

\def\eqbox#1{\tcboxmath[colback=banana,arc=0cm,frame hidden,enhanced]{#1}}
\def\figbox{\begin{center}\tcboxmath[colback=isabelline,frame hidden,enhanced,sharp corners]{\parbox[c][3cm][c]{10cm}{\ }}\end{center}}

\title[Quantum computing and quantum logic with trapped ions]{Quantum computing and quantum logic with trapped ions: Homework 3, due Nov.\ 4th, 2022}
\author{Tobias J.\ Osborne}
\date{\today}                                           % Activate to display a given date or no date

\begin{document}

\maketitle
Groups of two students can submit together. Calculations with computer algebra systems are acceptable if well documented (please also upload exported PDF). There will be a submission folder on Stud.IP. Please upload your solutions electronically and clearly state whom you are submiting for on the submission (name and student registration number)!

\setcounter{section}{1}

\begin{exercise}[1 mark]
   Show that the {\sc nand} gate can be used to simulate the {\sc and}, {\sc xor} and {\sc not} gates, provided wires and ancilla bits are available.
\end{exercise}

\begin{exercise}[1 mark]
  Consider a quantum computer with $3$ qubits, labelled $1$, $2$, and $3$. What is the matrix representation of the quantum circuit $U = \textsc{cnot}_{12}\textsc{cnot}_{23}$, where the subscripts indicate which qubits the gates act on. What is the quantum circuit diagram for this quantum circuit?
\end{exercise}

\begin{exercise}[3 marks]
  Argue that any classical circuit comprising {\sc and}, {\sc or}, etc.\ gates may be simulated reversibly, i.e., using invertible gates, using a quantum computer.
\end{exercise}

\begin{exercise}[1 mark]
  What is the matrix representation of a general controlled-$U$ operation, where $U$ is an arbitrary element of $\textsl{SU}(2)$, in the computational basis?
\end{exercise}

\begin{exercise}[3 marks]
  Find two-level unitary matrices $U_1$, $U_2$, $U_3$, so that
  \begin{equation}
    U_3 U_2 U_1 U = \mathbb{I},
  \end{equation}
  where $U$ is an arbitrary $2\times 2$ unitary matrix.
\end{exercise} 

\begin{exercise}[1 mark]
   Find a Gray code on $n=6$ bits connecting $\textbf{s} = 101001$ and $\textbf{t}=110011$.
\end{exercise}

\begin{exercise}[5 marks]
  Show how to effect the following two-level operation on $n=3$ qubits using controlled unitaries.
  \begin{equation}
    U = \begin{pmatrix}
      a & 0 & 0 & 0 & 0 & 0 & b & 0 \\
      0 & 1 & 0 & 0 & 0 & 0 & 0 & 0 \\
      0 & 0 & 1 & 0 & 0 & 0 & 0 & 0 \\
      0 & 0 & 0 & 1 & 0 & 0 & 0 & 0 \\
      0 & 0 & 0 & 0 & 1 & 0 & 0 & 0 \\
      0 & 0 & 0 & 0 & 0 & 1 & 0 & 0 \\
      c & 0 & 0 & 0 & 0 & 0 & d & 0 \\
      0 & 0 & 0 & 0 & 0 & 0 & 0 & 1 
    \end{pmatrix}.
  \end{equation}
  Here $a$, $b$, $c$, and $d$ are complex numbers so that $\widetilde{U}=\left(\begin{smallmatrix} a & b\\ c & d \end{smallmatrix}\right)$ is a unitary matrix. The Gray code 
  \begin{equation}
  \begin{matrix}
    000 \\
    010 \\
    110 \\
  \end{matrix}
\end{equation}
may be helpful.
\end{exercise} 

\end{document}
